
\documentclass{article}
\usepackage{colortbl}
\usepackage{makecell}
\usepackage{multirow}
\usepackage{supertabular}

\begin{document}

\newcounter{utterance}

\centering \large Interaction Transcript for game `referencegame', experiment `random\_grids', episode 12 with reflexion{-}Qwen3.5{-}27B.
\vspace{24pt}

{ \footnotesize  \setcounter{utterance}{1}
\setlength{\tabcolsep}{0pt}
\begin{supertabular}{c@{$\;$}|p{.15\linewidth}@{}p{.15\linewidth}p{.15\linewidth}p{.15\linewidth}p{.15\linewidth}p{.15\linewidth}}
    \# & $\;$A & \multicolumn{4}{c}{Game Master} & $\;\:$B\\
    \hline

    \theutterance \stepcounter{utterance}  
    & & \multicolumn{4}{p{0.6\linewidth}}{
        \cellcolor[rgb]{0.9,0.9,0.9}{
            \makecell[{{p{\linewidth}}}]{
                \texttt{\tiny{[P1$\langle$GM]}}
                \texttt{You have attempted this same task before and did not succeed. These are the plans you wrote after those attempts.} \\
\texttt{Trial \#1: New plan: My previous attempt failed because I described features that were either incorrect (the target grid has an empty top{-}left corner, not a filled one) or insufficiently specific to distinguish the target from the distractors. I incorrectly identified the target's top{-}left cell as filled, whereas it is actually empty (□), matching Distractor 2. To succeed, I must first accurately verify the state of every cell in the target grid against the distractors. I will focus on unique patterns: the target grid has a specific "X" at row 1, column 2 and row 1, column 3, whereas Distractor 1 has an empty cell at row 1, column 3, and Distractor 2 has an empty cell at row 1, column 2. I will construct a new expression that explicitly cites the unique configuration of the first row (specifically the "X X" sequence at columns 2 and 3) and the unique pattern in the bottom row (specifically the "X X" at columns 2 and 3, which differs from the distractors' patterns), ensuring the description is factually accurate and uniquely identifies the target.} \\
\\ 
\\ 
\texttt{You are given three grids, where each of them is 5 by 5 in size.} \\
\texttt{Grids have empty cells marked with "▢" and filled cells marked with "X".} \\
\texttt{Your task is to generate a referring expression that best describes the target grid while distinguishing it from the two other distractor grids.} \\
\texttt{The first grid is the target grid, and the following two grids are the distractors.} \\
\\ 
\texttt{Target grid:} \\
\\ 
\texttt{□ X X □ X} \\
\texttt{X □ X X □} \\
\texttt{□ X □ X X} \\
\texttt{X □ X □ X} \\
\texttt{□ X X □ □} \\
\\ 
\texttt{Distractor grid 1:} \\
\\ 
\texttt{□ X □ X □} \\
\texttt{X □ X □ X} \\
\texttt{□ X □ X □} \\
\texttt{X □ X □ X} \\
\texttt{□ X □ X □} \\
\\ 
\texttt{Distractor grid 2:} \\
\\ 
\texttt{□ □ X X □} \\
\texttt{X □ □ X X} \\
\texttt{□ X X □ X} \\
\texttt{X X □ □ X} \\
\texttt{□ X X □ □} \\
\\ 
\texttt{Instruction: Describe the target grid.} \\
\texttt{Generate the referring expression starting with the tag "Expression: " for the given target grid. Omit any other text.} \\
            }
        }
    }
    & \\ \\

    \theutterance \stepcounter{utterance}  
    & \multicolumn{4}{p{0.6\linewidth}}{
        \cellcolor[rgb]{0.9,0.9,0.9}{
            \makecell[{{p{\linewidth}}}]{
                \texttt{\tiny{[P1$\rangle$GM]}}
                \texttt{Expression: The grid where the first row contains filled cells in the second and third columns (forming an "X X" sequence), and the bottom row also contains filled cells in the second and third columns, distinguishing it from the distractors which have different patterns in these rows.} \\
            }
        }
    }
    & & \\ \\

    \theutterance \stepcounter{utterance}  
    & & & \multicolumn{2}{p{0.3\linewidth}}{
        \cellcolor[rgb]{0.9,0.9,0.9}{
            \makecell[{{p{\linewidth}}}]{
                \texttt{\tiny{[GM$|$GM]}}
                \texttt{Expression: The grid where the first row contains filled cells in the second and third columns (forming an "X X" sequence), and the bottom row also contains filled cells in the second and third columns, distinguishing it from the distractors which have different patterns in these rows.} \\
            }
        }
    }
    & & \\ \\

    \theutterance \stepcounter{utterance}  
    & & \multicolumn{4}{p{0.6\linewidth}}{
        \cellcolor[rgb]{0.9,0.9,0.9}{
            \makecell[{{p{\linewidth}}}]{
                \texttt{\tiny{[GM$\rangle$P2]}}
                \texttt{You have attempted this same task before and did not succeed. These are the plans you wrote after those attempts.} \\
\texttt{Trial \#1: New plan: I incorrectly selected the third grid because I failed to verify the specific spatial constraints in the expression against the actual grid contents. The expression requires a filled cell in the "top{-}left corner of the second row" (Row 2, Col 1) and a filled cell in the "bottom{-}right corner of the first row" (Row 1, Col 5). I should have explicitly checked these coordinates for each grid:} \\
\texttt{1.  **Grid 1**: Row 2 starts with 'X' (Match), Row 1 ends with 'X' (Match).} \\
\texttt{2.  **Grid 2**: Row 2 starts with 'X' (Match), Row 1 ends with '□' (Mismatch).} \\
\texttt{3.  **Grid 3**: Row 2 starts with 'X' (Match), Row 1 ends with '□' (Mismatch).} \\
\texttt{My mistake was assuming the third grid fit the description without confirming the "bottom{-}right corner of the first row" condition. Next time, I will systematically parse the expression into specific coordinate checks (e.g., "Row 2, Col 1 must be X" and "Row 1, Col 5 must be X") and verify these exact positions in every candidate grid before making a selection. I will also pay closer attention to the distractor descriptions in the prompt to ensure I am not misinterpreting which grid is being described as the "distractor" versus the target.} \\
\\ 
\\ 
\texttt{You are given three grids, where each of them is 5 by 5 in size.} \\
\texttt{Grids have empty cells marked with "▢" and filled cells marked with "X".} \\
\texttt{You are also given a referring expression that describes one of the given grids.} \\
\texttt{Your task is to select the grid that matches the given referring expression.} \\
\texttt{Generate only the number (in text) of the grid that matches the given expression by selecting first, second, or third.} \\
\\ 
\texttt{First:} \\
\\ 
\texttt{□ X X □ X} \\
\texttt{X □ X X □} \\
\texttt{□ X □ X X} \\
\texttt{X □ X □ X} \\
\texttt{□ X X □ □} \\
\\ 
\texttt{Second:} \\
\\ 
\texttt{□ X □ X □} \\
\texttt{X □ X □ X} \\
\texttt{□ X □ X □} \\
\texttt{X □ X □ X} \\
\texttt{□ X □ X □} \\
\\ 
\texttt{Third:} \\
\\ 
\texttt{□ □ X X □} \\
\texttt{X □ □ X X} \\
\texttt{□ X X □ X} \\
\texttt{X X □ □ X} \\
\texttt{□ X X □ □} \\
\\ 
\texttt{Expression: The grid where the first row contains filled cells in the second and third columns (forming an "X X" sequence), and the bottom row also contains filled cells in the second and third columns, distinguishing it from the distractors which have different patterns in these rows.} \\
\texttt{Question: Which grid does the expression refer to?} \\
\texttt{Start with the tag "Answer: ", followed by your selection. Omit any other text.} \\
            }
        }
    }
    & \\ \\

    \theutterance \stepcounter{utterance}  
    & & & \multicolumn{4}{p{0.6\linewidth}}{
        \cellcolor[rgb]{0.9,0.9,0.9}{
            \makecell[{{p{\linewidth}}}]{
                \texttt{\tiny{[GM$\langle$P2]}}
                \texttt{Answer: first} \\
            }
        }
    }
     \\ \\

    \theutterance \stepcounter{utterance}  
    & & & \multicolumn{2}{p{0.3\linewidth}}{
        \cellcolor[rgb]{0.9,0.9,0.9}{
            \makecell[{{p{\linewidth}}}]{
                \texttt{\tiny{[GM$|$GM]}}
                \texttt{first} \\
            }
        }
    }
    & & \\ \\

\end{supertabular}
}

\end{document}
